% Modernised reading — fonds Alexandre Grothendieck, shelfmark 86 (pages 1-32,
% the whole folder).
%
% This is an INTERPRETATION, not a transcription. It re-reads the pages in
% today's notation and vocabulary, states the hypotheses the manuscript leaves
% implicit, and drops the critical apparatus so that the mathematics reads
% straight through. Everything the brackets and underlines carried moves into a
% footnote. It is held to being mathematically correct as it stands; where the
% manuscript is loose or mistaken, the true statement is what appears here and
% the footnote says what the page has. In any dispute of substance the
% transcription governs: whoever cites Grothendieck must cite batch-01.fr.tex
% (pp. 1-20) or batch-02.fr.tex (pp. 21-32).
%
% Pass: Opus 5.5 (claude-opus-5-5), 2026-10-10 — first pass, unchecked against the pages by a human.
% The folder was read whole, its two batches together; this is one document
% for the shelfmark. It was made on Opus 5.5 and has not been measured against
% the reference reading of folder 115 (made on Opus 5). The literature named in
% the footnotes is cited from memory and was not checked against the sources.
% Every count and every identity on the icosahedron asserted below (rho^3 = id,
% the five trièdres, the fixed trièdres of a face, the edge census of (C), the
% dictionary of pp. 28-30, the permutations of pp. 31-32) was checked by
% machine on the regular icosahedron with golden-ratio coordinates.
%
% Scope. Pages 1, 9 and 17 are covers whose words head the three parts here
% (« Polyèdres réguliers I », « II » with « géométrie d'incidence et ens. à
% opérateurs », « Cours / Polyèdres réguliers III »); no \pagerange claims
% them. Pages 31-32 are in another ink.
%
% Spine.
%   pp. 2-4    I: combinatorial polyhedra as ordered sets S u A u F, axioms
%              1)-4); 4) is the lattice condition; Phi -> P(S) an order
%              embedding; every face has >= 3 sides.
%   p. 5       plan of a course on the icosahedron.
%   pp. 6-8    three descriptions of the icosahedron adapted to an axis through
%              two opposite vertices, faces, edges; stabilisers of order 20,
%              12, 8.
%   pp. 10-16  II: flags, Rep, the involutions sigma_0, sigma_1, sigma_2; the
%              functor to sets under the cartographic group is fully faithful;
%              connectedness, regularity, stabilisers D_p, D_q, (Z/2)^2.
%   pp. 18-25  III: the antipodism iota, the 15 edge directions, rho of order
%              3, the five « trièdres distingués » T; P(s) = T; G -> Alt_T onto
%              with kernel <iota>; G = Z/2 x G+, G+ = A_5.
%   pp. 26-30  reconstruction of the « icosaèdre gauche » Pi/iota (today's
%              hemi-icosahedron) from the oriented five-element set (T, varpi):
%              sets, incidences, flags = numberings compatible with varpi.
%   pp. 31-32  other ink: table of the flag-type sets of the icosaèdre gauche
%              over an oriented five-element set E; sigma_i as permutations.
%
% Notation fixed for the document. Phi = S u A u F ordered; Rep(Pi) = E the set
% of flags (his « repères »); hat-Gamma the group <sigma_0, sigma_1, sigma_2 |
% sigma_i^2, (sigma_0 sigma_2)^2>, which the pages write with a gothic S_2
% (not the symmetric group) and the reading of folder 88 writes hat-Gamma;
% G = Aut(Pi); bar-sigma_i the automorphisms of p. 14. D_n is dihedral of
% order 2n. The antipodism, an underlined a on the pages, is iota here, to keep
% a for edges; the trièdres and the map tr keep his names. On pp. 7-8 his T,
% T_1, T'_1, T' (the layers of (B)) are Delta, Delta_1, Delta'_1, Delta', and
% his sigma, sigma' (the two opposite edges of (C)) are e, e', to keep T and
% sigma free. « Orientation » of a finite set X (card >= 2) is an
% Alt_X-orbit of numberings; for a polygon it is NOT the choice of a cyclic
% direction of folder 89 — stated in the spine section.
%
% Substantive departures from the pages, each footnoted where it occurs:
%  - p. 2: axiom 4) and its unpacking a)-b), the NB and the box (b) are
%    replaced by the clean equivalent (i) A -> P_2(S) injective and (ii) two
%    distinct faces meet in nothing, a vertex, or an edge with its two
%    vertices; proved here. « sup majorés <=> inf minorés » is true for finite
%    ordered sets; said so.
%  - pp. 2-4: Cor. 1 (p. 2) is proved from Cor. Main 2 (p. 4); the order of
%    proof is reversed.
%  - p. 6: « delta(t, s') = 1 » is 2 when s -> s' is the antipodism (the only
%    reading under which the parenthesis = delta(t', s) holds); no bijection
%    makes 1 true.
%  - p. 8: the count 4 for the edges {Pi_{s,u}, Pi_{s,v}} is 2 (the primed
%    pairs are the same edges); with it the census sums to 30.
%  - p. 14: g -> bar-g is an ANTI-homomorphism (the page: « hom. », rest
%    illegible).
%  - p. 15: 3 <= p, q needs axiom 4) (Cor. 3), which a mere incidence geometry
%    does not have.
%  - p. 23: the first triple of directions read {u,v}, {v,w_1}, {w,u_1} is
%    {u,v_1}, {v,w_1}, {w,u_1}: {u,v} is an edge of f.
%  - p. 27: which packet of A(T) is called A+(T, varpi) is a convention the
%    page fixes by notation only.
%  - p. 28: b) completed (palimpsest on the page): a and f incident iff
%    t not in Omega and Omega is a side of gamma not in alpha. Ours.
%  - p. 31: A^up is the set of POINTED triangles (30), not Tr_3(E) (10); the
%    page's own comment says « pointés ». The sigma_1 of p. 31 is not the one
%    the dictionary of p. 30 gives ((a5,a4,a3,a2,a1)); both triples satisfy
%    the relations and give isomorphic flag actions.
%  - p. 32: « sigma = sigma_0 sigma_1 » has the value of sigma_0 sigma_2.
%
% Announced and never established in the folder: the proofs of Cor. Main 2
% for faces (p. 4, « … ») and of Corollaire 4 (p. 4, mostly illegible); the
% list of the regular polyhedra realisable as convex polyhedra (p. 5, breaks
% off); the « interprétation concrète en termes de coloriage » (p. 19, breaks
% off); the description of S~, A~, F~ in terms of (T, varpi) and of Pi-bar as
% a quotient polyhedron (p. 30, « Il reste »), of which the table of p. 31
% gives the sets but no incidences and no proof.
\documentclass[11pt,a4paper]{article}
\input{../preamble/grothendieck}

\folder{86}
\pages{1}{32}
\dating{[à partir de 1977]}
\watermark{Édition de démonstration\\interprétation personnelle de l'œuvre}
\foldertitle{[Polyèdres réguliers] : notes manuscrites (s.d.) — lecture modernisée du dossier entier}

\begin{document}

\section*{Résumé}

\begin{resume}
L'icosaèdre est le solide de Platon à vingt faces triangulaires, douze sommets
et trente arêtes. Ces trente-deux pages se demandent ce qu'il en reste quand on
oublie tout de la géométrie --- les longueurs, les angles, l'espace où il est
posé --- et qu'on ne garde que la liste de ses sommets, de ses arêtes, de ses
faces, et de qui touche qui. La réponse est : à peu près tout, et même un peu
plus, car on y voit alors une structure que la géométrie cache.

Le premier problème est de dire ce qu'est un polyèdre sans coordonnées. On ne
peut pas parler de « solide convexe » ; on dispose seulement d'une relation
d'incidence. Grothendieck écrit les axiomes qui font d'une telle relation un
polyèdre : chaque arête a deux bouts et borde deux faces, autour d'une face on
fait un tour complet, et deux faces ne se touchent que le long d'une arête ou
en un sommet. Puis il change de point de vue. Au lieu de compter les sommets,
arêtes et faces, il compte les \emph{drapeaux} : un sommet, une arête qui en
part, une face que borde cette arête. Sur les drapeaux agissent trois gestes
élémentaires --- changer le sommet, changer l'arête, changer la face, en
gardant les deux autres --- et ces trois gestes suffisent à reconstituer le
polyèdre entier. C'est comme décrire une ville non par son plan, mais par
l'ensemble des positions d'un piéton et les trois pas qu'il peut faire.

L'idée centrale arrive dans la troisième partie, intitulée « Cours ». Les
trente arêtes de l'icosaèdre vont par paires parallèles : il y a quinze
directions. Grothendieck définit, par une construction purement combinatoire,
une permutation de ces quinze directions dont le cube est l'identité ; elle
les regroupe en cinq « trièdres distingués », trois directions chacun. Dans
l'icosaèdre réel, ce sont cinq repères orthogonaux de l'espace. Toute
symétrie de l'icosaèdre permute ces cinq trièdres, et l'on obtient ainsi,
sans un seul calcul de coordonnées, le fait classique que les rotations de
l'icosaèdre forment le groupe des permutations paires de cinq objets. Il va
plus loin : il reconstitue l'icosaèdre « replié » sur lui-même par la symétrie
centrale --- un polyèdre à six sommets, quinze arêtes et dix faces, qui ne
vit plus dans l'espace mais sur le plan projectif --- à partir des seuls cinq
trièdres et d'un choix de sens : les sommets deviennent des pentagones dessinés
sur cinq points, les faces des paires de points, les drapeaux des façons de
numéroter les cinq points.

Ce qu'on y voit de sa manière : l'objet le plus classique est repris comme s'il
était inconnu, et chaque construction est faite pour ne dépendre d'aucun choix.
Les trièdres sont définis sans parler d'angle droit, et l'orientation de
l'ensemble des cinq trièdres n'est pas choisie : elle est lue sur l'icosaèdre,
où elle était déjà. On cherchera ensuite, sous les noms d'aujourd'hui, les
polytopes abstraits réguliers, le groupe cartographique, l'hémi-icosaèdre, et
l'isomorphisme entre le groupe de l'icosaèdre et le groupe alterné $A_5$.

\keywords{abstract polytope, flag, lattice, incidence geometry, regular polyhedron, icosahedron, cartographic group, alternating group A5, compound of five octahedra, 1-factorization, hemi-icosahedron, orientation of a finite set}
\end{resume}

\pagerange{2}{32}
\section*{Le fil du dossier, et les conventions}

Trois chemises (pages 1, 9 et 17) portent les titres de sa main,
« Polyèdres réguliers I », « II » et « Cours --- Polyèdres réguliers
III » ; le titre de l'inventaire en vient vraisemblablement. Il n'y a pas de
pagination de sa main ; la page 10 renvoie à une « p.~1 » qui est
vraisemblablement la page 2.

\begin{itemize}
  \item \textbf{Pages 2 à 4 (I) : polyèdres combinatoires.} Un ensemble ordonné
    $S \sqcup A \sqcup F$, quatre axiomes ; le quatrième est une condition de
    treillis. Toute facette est la borne supérieure de ses sommets ; toute face
    a au moins trois côtés.
  \item \textbf{Page 5 : un plan de cours} sur l'icosaèdre.
  \item \textbf{Pages 6 à 8 : trois descriptions de l'icosaèdre}, adaptées à
    un axe passant par deux sommets, deux faces, deux arêtes opposés.
  \item \textbf{Pages 10 à 16 (II) : géométries d'incidence et ensembles à
    opérateurs.} Drapeaux, trois involutions, un théorème de pleine fidélité ;
    connexité, régularité, stabilisateurs.
  \item \textbf{Pages 18 à 25 (III) : les trièdres distingués et le groupe
    alterné.} La permutation $\rho$, les cinq trièdres, l'isomorphisme
    $G^{+} \simeq \operatorname{Alt}_T$.
  \item \textbf{Pages 26 à 30 : reconstruction de l'icosaèdre gauche} à partir
    de l'ensemble orienté des cinq trièdres.
  \item \textbf{Pages 31 et 32}, d'une autre encre : un tableau des ensembles
    de drapeaux de l'icosaèdre gauche, et les trois involutions écrites comme
    permutations de cinq lettres.
\end{itemize}

\textbf{Les dossiers voisins.} La partie II est le cadre dans lequel la
lecture du dossier 88 (« Cartes standards ») est écrite : les mêmes
involutions $\sigma_0, \sigma_1, \sigma_2$ sur les repères, le même groupe,
noté là $\widehat{\Gamma}$, les mêmes rotations $\sigma_1\sigma_2$ autour d'un
sommet et $\sigma_0\sigma_1$ autour d'une face, et l'icosaèdre y est la carte
standard $C_{5,3}$. Les polygones combinatoires dont les pages 2 et 15 font les
« entourages » des sommets et des faces sont ceux des pages 19 et 20 du
dossier 89.

\textbf{Conventions, valables pour tout le dossier.} $\Phi = S \sqcup A \sqcup
F$ est l'ensemble ordonné des \emph{facettes} (sommets, arêtes, faces),
$E = \operatorname{Rep}(\Pi)$ l'ensemble des drapeaux maximaux --- ses
« repères » ---, et
\[
\widehat{\Gamma} = \bigl\langle \sigma_0, \sigma_1, \sigma_2 \bigm|
\sigma_0^2 = \sigma_1^2 = \sigma_2^2 = (\sigma_0\sigma_2)^2 = 1 \bigr\rangle ,
\]
qu'il note d'un grand $\mathfrak{S}$ gothique d'indice $2$, qui n'est pas le
groupe symétrique\footnote{On adopte la lettre de la lecture du dossier 88,
où le groupe porte ce nom sur la page. Le nom de \emph{groupe cartographique}
est celui de l'\emph{Esquisse d'un programme} (1984) ; il n'est pas dans ce
dossier.}. $G = \operatorname{Aut}(\Pi)$ ; $D_n$ est le groupe diédral
d'ordre $2n$. L'antipodisme de l'icosaèdre, qu'il note d'un $a$ souligné, est
noté ici $\iota$, pour laisser la lettre $a$ aux arêtes.

\textbf{Deux sens du mot « orientation ».} Dans la troisième partie, une
\emph{orientation} d'un ensemble fini $X$ de cardinal $\geq 2$ est une orbite
du groupe alterné $\operatorname{Alt}_X$ sur les numérotations de $X$ ; il y en
a deux. Pour un polygone, ce n'est pas la même chose que l'orientation des
pages 19 et 20 du dossier 89, qui est le choix de l'un des deux sens de
parcours : les deux sens de parcours d'un pentagone définissent la
\emph{même} orientation de l'ensemble de ses sommets (page 21). Pour un
triangle, les deux notions coïncident, et la page 23 écrit
$\operatorname{Or}(f)$ pour l'ensemble des deux ordres circulaires.

\pagerange{2}{4}
\section*{I. Polyèdres combinatoires (pages 2 à 4)}

\pagerange{2}{2}
\subsection*{Les axiomes}

On se donne des ensembles finis $S$, $A$, $F$ (sommets, arêtes, faces) et des
relations d'incidence $R \subset S \times A$, $R' \subset A \times F$. Elles
font de $\Phi = S \sqcup A \sqcup F$ un ensemble ordonné, où $s < a$ signifie
$(s, a) \in R$, $a < f$ signifie $(a, f) \in R'$, et $s < f$ signifie qu'il
existe $a$ avec $s < a < f$ : $R'' = R \circ R'$. Les axiomes sont :
\begin{enumerate}
  \item[1)] toute arête a exactement deux sommets et est bordée par exactement
    deux faces ;
  \item[2)] (propriété du losange) si $s < f$, il y a exactement deux arêtes
    $a$ avec $s < a < f$ ;
  \item[3)] pour toute face $f$, l'ensemble ordonné $\Phi_{<f}$ est connexe ;
  \item[3')] pour tout sommet $s$, l'ensemble ordonné $\Phi_{>s}$ est connexe ;
  \item[4)] dans $\Phi$, deux éléments qui ont un majorant commun ont une borne
    supérieure.
\end{enumerate}
Par 1) et 2), $\Phi_{<f}$ est un graphe biparti où chaque sommet et chaque arête
sont de degré $2$, c'est-à-dire une réunion disjointe de polygones ; 3) dit
que c'est un seul polygone combinatoire, à $q(f) \geq 2$ côtés. De même,
$\Phi_{>s}$ est un polygone à $p(s) \geq 2$ côtés, dont les « sommets » sont
les arêtes issues de $s$ et les « côtés » les faces qui contiennent
$s$\footnote{Ces polygones sont ceux des pages 19 et 20 du dossier 89, sous
sa première description (un ensemble ordonné de sommets et d'arêtes, chacun
de degré $2$). La page qualifie $\Phi_{<f}$ de « contour ». Une note verticale,
le long de l'accolade qui embrasse 1) à 3'), n'est pas lue, sauf peut-être son
dernier mot, « connexité ».}. Aujourd'hui, 1) à 3') définissent un
\emph{polyèdre abstrait} --- un polytope abstrait de rang $3$, au sens par
exemple du livre de McMullen et Schulte (2002) ---, et 4) demande que
$\Phi$, augmenté d'un plus petit et d'un plus grand élément, soit un
treillis\footnote{Les noms ne sont pas sur la page. La page écrit les axiomes
pour un polyèdre fini.}.

\textbf{L'axiome 4), symétrique et explicité.} Pour un ensemble ordonné fini,
4) équivaut à l'énoncé dual, « deux éléments qui ont un minorant commun ont
une borne inférieure » : l'un et l'autre disent que $\widehat{\Phi} = \Phi \sqcup
\{\hat 0, \hat 1\}$ est un treillis, et un ensemble ordonné fini qui a toutes
les bornes supérieures de deux éléments et un plus petit élément a aussi toutes
les bornes inférieures. Explicitement, \emph{sous 1) à 3'), l'axiome 4)
équivaut à la conjonction de}
\begin{enumerate}
  \item[(i)] une arête est déterminée par ses deux sommets : $A \to
    \mathfrak{P}_2(S)$ est injective ;
  \item[(ii)] deux faces distinctes ont en commun, comme facettes inférieures,
    soit rien, soit un sommet, soit une arête et ses deux sommets.
\end{enumerate}
En effet, deux éléments distincts qui ont un minorant commun sont deux arêtes,
une arête et une face, ou deux faces, et leurs minorants communs sont leurs
sommets communs (et, pour deux faces, leurs arêtes communes). La borne
inférieure existe si et seulement si cet ensemble a un plus grand élément :
pour deux arêtes, il faut qu'elles aient au plus un sommet commun, ce qui est
(i) ; pour deux faces, c'est (ii). Reste le cas d'une arête $a = \{s, s'\}$ et
d'une face $f$ qui ne la contient pas, dont les sommets communs doivent être
au plus un : s'ils étaient $s$ et $s'$, la face $f'$ qui borde $a$ rencontrerait
$f$ en $s, s'$, donc, par (ii), en une arête de sommets $s, s'$, distincte de
$a$, contre (i)\footnote{La page procède autrement. Elle déplie 4) en a) « une
arête est déterminée par l'ensemble de ses extrémités, et les faces qui
majorent [celles-ci] majorent l'arête » et b) « une face est connue par une
arête et un sommet non incident, [ou] par deux sommets non incidents à une
même arête », lecture en partie incertaine. Un NB, aux alinéas numérotés 1°)
à 3°) (deux fois 3°)), analyse correctement, cas par cas, quand
$\operatorname{Sup}(s, s')$, $\operatorname{Sup}(s, a)$,
$\operatorname{Sup}(a, a')$ existent : c'est quand le majorant évident est le
seul. Un encadré en tire la forme (ii), suivie de « $A \hookrightarrow
\mathfrak{P}_2(S)$ », l'ordre des deux membres étant incertain ; une note
oblique de la marge en écrit le dual, $A \hookrightarrow \mathfrak{P}_2(F)$,
et un 4') dual de a)-b). L'équivalence telle qu'on l'énonce, et sa preuve,
sont de l'édition.}.

\pagerange{2}{4}
\subsection*{Les premières conséquences}

\textbf{Corollaire 3} (page 4). \emph{Toute face a au moins trois côtés} ;
dualement, \emph{de tout sommet partent au moins trois arêtes} (page 3). Une
face à deux côtés aurait deux arêtes de mêmes extrémités, contre (i).

\textbf{Corollaire 2} (page 4, « Cor.~Main 2 »). \emph{Toute facette est la
borne supérieure des sommets qu'elle majore.} C'est clair pour un sommet. Pour
une arête $a = \{s, s'\}$ : $a$ majore $s$ et $s'$, et tout autre majorant
commun est une face qui contient $s$ et $s'$, donc l'arête de sommets $s, s'$
par (ii), c'est-à-dire $a$ par (i) ; donc $a = \operatorname{Sup}(s,s')$. Pour
une face $f$ : la borne supérieure de deux arêtes distinctes de $f$ est une
face majorée par $f$, donc $f$, et, $f$ ayant au moins deux arêtes\footnote{La page écrit cette chaîne d'égalités, « cqfd », avec « cor.~2 »
au-dessus de la deuxième, après un début de justification en partie illisible
qui invoque la connexité ; la preuve pour une face, page 4, s'arrête sur
« pour face cela résulte de … ». Les justifications sont de l'édition.},
\[
\operatorname*{Sup}_{s < f} s \;=\; \operatorname*{Sup}_{a<f}\,
\operatorname*{Sup}_{s<a} s \;=\; \operatorname*{Sup}_{a<f} a \;=\; f
\]
.

\textbf{Corollaire 1} (pages 2 à 4). \emph{L'application $\alpha : \Phi \to
\mathfrak{P}(S)$, qui envoie une facette sur l'ensemble des sommets qu'elle
majore, est un plongement d'ensembles ordonnés : $\alpha(x) \subset \alpha(y)$
si et seulement si $x \leq y$. Son image, augmentée de $\varnothing$, est
stable par intersection.} Si $\alpha(x) \subset \alpha(y)$, $y$ majore tous
les sommets de $x$, donc leur borne supérieure, qui est $x$ par le corollaire 2.
Si $x$ et $y$ ont un minorant commun, ils ont un sommet commun, et
$\alpha(x) \cap \alpha(y) = \alpha(\operatorname{Inf}(x, y))$ ; sinon
$\alpha(x) \cap \alpha(y) = \varnothing$\footnote{La page énonce ce corollaire
page 2, avant le corollaire 2 dont sa preuve dépend, et l'achève page 4 ; la
fin de la démonstration de a) est en blanc (« résulte dans … »). On a remis
la preuve dans l'ordre logique.}. Dualement (Corollaire 4, page 4), un sommet
est la borne inférieure des faces qui le contiennent, et $S \to
\mathfrak{P}(F)$ est injective\footnote{Le corollaire 4 est en grande partie
illisible ; on en donne l'énoncé que dicte le mot « Dualement » qui le
précède, et les mots lus : « des faces qui l'entourent », « une injection
$S \to \mathfrak{P}(F)$ ».}.

\textbf{Remarque} (page 3). Pour l'icosaèdre, tout se lit sur le graphe
$(S, A)$, $A \subset \mathfrak{P}_2(S)$ : les faces sont exactement les
triangles du graphe, $F = \{f \in \mathfrak{P}_3(S) \mid \mathfrak{P}_2(f)
\subset A\}$, et $\operatorname{Aut}(\Pi)$ est le groupe des permutations de
$S$ qui envoient arêtes sur arêtes, c'est-à-dire qui conservent la distance
combinatoire\footnote{Que tout triangle du graphe soit une face tient à ce que
les voisins d'un sommet forment un pentagone, sans corde : il y a $12 \cdot 5
/ 3 = 20$ triangles. Justification de l'édition. Une note en diagonale dans
la marge n'est pas lue.}.

\pagerange{5}{5}
\section*{Un plan de cours sur l'icosaèdre (page 5)}

La page 5 est un plan, sous le titre « Icosaèdre ». 1. Une description concrète,
par un modèle ou un dessin. 2. Le dénombrement : $12$ sommets, $30$ arêtes,
$20$ faces, $p = 5$ arêtes en chaque sommet, $q = 3$ côtés à chaque face, et
$2 \cdot 30 = 5 \cdot 12 = 3 \cdot 20$. Les propriétés polyédrales, comme
géométrie d'incidence et comme polyèdre combinatoire. 3. Les distances
combinatoires et l'existence des bornes\footnote{La ligne 3 est très
surchargée ; seuls « Distances combinatoires » et « Sup et Inf » sont sûrs.}.
4. Les « perspectives », dont la première, par un couple de sommets opposés,
a pour groupe de symétrie $D_5$ ou $\mathbf{Z}/2 \times D_5$ ; une parenthèse
sur les groupes d'automorphismes des polygones combinatoires --- les groupes
diédraux. Puis la transitivité de $G$ sur les sommets, sur les couples
$s < a$, sur les triplets $s < a < f$ : l'icosaèdre est un \emph{polyèdre
combinatoire régulier}. Un NB ajoute qu'il y en a une infinité d'autres, mais
que parmi eux seuls quelques-uns sont réalisables comme polyèdres convexes de
l'espace de dimension $3$ ; une note marginale propose d'interpréter
l'antipodisme comme automorphisme combinatoire\footnote{Après « seuls » la page
est illisible. Les polyèdres réguliers convexes sont les cinq solides de
Platon ; ce nombre n'est pas lu sur la page.}.

\pagerange{6}{8}
\section*{Trois descriptions de l'icosaèdre (pages 6 à 8)}

Chacune est adaptée à un axe passant par deux éléments opposés ; elle donne
l'ensemble des sommets, les arêtes par types avec leur nombre, et le groupe des
symétries qui conservent l'axe, d'ordre $120$ divisé par le nombre d'axes du
type. Les figures qui les accompagnent --- vues le long de l'axe, projections,
un icosaèdre inscrit dans un cube --- ne sont pas reproduites.

\pagerange{6}{6}
\subsection*{(A) Par deux sommets opposés}

$S = \{o\} \sqcup \Pi \sqcup \Pi' \sqcup \{o'\}$, où $\Pi$ et $\Pi'$ sont les
pentagones des voisins de $o$ et de $o'$, et $s \mapsto s'$ l'antipodisme, qui
échange $\Pi$ et $\Pi'$. Les arêtes sont les $\{o, s\}$, $s \in \Pi$, et
$\{o', s'\}$ ($10$) ; les côtés des deux pentagones ($10$) ; et les $\{s, t\}$,
$s \in \Pi$, $t \in \Pi'$, avec $\delta(t, s') = 2$, $\delta$ étant la distance
combinatoire ($10$) : chaque sommet d'un pentagone est relié aux deux sommets de
l'autre qui ne sont ni l'antipode du sien ni ses voisins\footnote{La page écrit
« $\delta(t, s') = 1$ ($= \delta(t', s)$) », la lettre $s'$ étant d'une
lecture incertaine. Avec $\delta(t,s') = \delta(t',s)$, que donne une isométrie
involutive comme l'antipodisme, la valeur est $2$ : dans l'icosaèdre,
$\delta(x, \iota y) = 3 - \delta(x, y)$. Aucune bijection $\Pi \simeq \Pi'$ ne
rend la valeur $1$ correcte, car les deux voisins de $s$ dans $\Pi'$ sont
adjacents entre eux, et les deux sommets à distance $1$ d'un même sommet ne le
sont pas.}. Groupe de symétrie $\mathbf{Z}/2 \times D_5$, d'ordre $20 = 120/6$,
$6 = 12/2$ étant le nombre d'axes par deux sommets.

\pagerange{7}{8}
\subsection*{(B) Par deux faces opposées}

$S = \Delta \sqcup \Delta_1 \sqcup \Delta'_1 \sqcup \Delta'$ : deux faces
opposées $\Delta$ et $\Delta'$, et deux couches de trois sommets, avec des
bijections $t \mapsto t_1 \mapsto t'_1$ et $t \mapsto t'$, où $t_1$ est le
troisième sommet de la face qui borde le côté de $\Delta$ opposé à
$t$\footnote{La page note les quatre couches $T$, $T_1$, $T'_1$, $T'$ ; on
garde $T$ pour les trièdres de la troisième partie. La règle qui fixe $t_1$ est
de l'édition ; la page écrit les bijections dans un tableau.}. Arêtes :
$\mathfrak{P}_2(\Delta)$ et $\mathfrak{P}_2(\Delta')$ ($6$) ; $\{t, u\}$,
$t \in \Delta$, $u \in \Delta_1$, $u \neq t_1$, et de même en primant ($12$) ;
$\{t_1, u\}$, $u \in \Delta'_1$, $u \neq t'_1$ ($6$) ; $\{t, t'_1\}$ et
$\{t', t_1\}$ ($6$) ; au total $30$\footnote{Les indices et accents des deux
dernières lignes sont d'une lecture incertaine ; la lecture retenue est celle
qui donne à chaque sommet ses cinq voisins.}. Groupe de symétrie
$\mathbf{Z}/2 \times \mathfrak{S}_3$, d'ordre $12 = 120/10$, $10 = 20/2$.

\pagerange{8}{8}
\subsection*{(C) Par deux arêtes opposées}

Soit $e = \{s, t\}$ une arête, $\varphi = \{u, v\}$ les troisièmes sommets des
deux faces qui la bordent, et $e'$, $\varphi'$ leurs antipodes. Pour $x \in e$,
$y \in \varphi$, notons $\Pi_{x,y}$ le voisin de $x$ qui est adjacent à $y$ et
distinct de l'autre extrémité de $e$. Les quatre sommets restants sont les
$\Pi_{x,y}$, et la même construction faite à partir de $e'$ donne les mêmes :
$\Pi_{s',u'} = \Pi_{t,v}$, $\Pi_{s',v'} = \Pi_{t,u}$, etc. D'où
\[
S = e \sqcup \varphi \sqcup e' \sqcup \varphi' \sqcup \bigl[(e \times \varphi)
\sqcup (e' \times \varphi')\bigr] / (\mathbf{Z}/2) ,
\]
où l'élément non trivial de $\mathbf{Z}/2$ envoie $(x, y)$ sur
$((\varepsilon x)', (\varepsilon y)')$, $\varepsilon$ échangeant les deux
éléments de $e$ et de $\varphi$ ; l'antipodisme agit sur les $\Pi_{x,y}$ par
$(x, y) \mapsto (\varepsilon x, \varepsilon y)$\footnote{La page note les deux
arêtes $\sigma$, $\sigma'$, et l'élément de $\mathbf{Z}/2$ encore $\sigma$ ; on
les renomme. La ligne est surchargée, une première identification est biffée,
et le $\mathbf{Z}/2\mathbf{Z}$ est d'une lecture incertaine ; la description
qu'on donne est celle que confirment les égalités de la figure
($\Pi_{t,v} = \Pi_{s',u'}$, $\Pi_{s,u} = \Pi_{t',v'}$, …).}. Les arêtes :
\[
\begin{array}{lr}
e,\ e' & 2 \\
\{x, y\},\ \{x', y'\} \quad (x \in e,\ y \in \varphi) & 8 \\
\{\Pi_{s,u}, \Pi_{s,v}\},\ \{\Pi_{t,u}, \Pi_{t,v}\} & 2 \\
\{y, \Pi_{x,y}\},\ \{y', \Pi_{x',y'}\} & 8 \\
\{x, \Pi_{x,y}\},\ \{x', \Pi_{x',y'}\} & 8 \\
\{u, v'\},\ \{v, u'\} & 2
\end{array}
\]
au total $30$\footnote{La page compte $4$ pour la troisième ligne, et le total
fait alors $32$ : les paires $\{\Pi_{s',u'}, \Pi_{s',v'}\}$ et
$\{\Pi_{t',u'}, \Pi_{t',v'}\}$ sont les deux arêtes déjà comptées. Le « $2$ »
de la première ligne est d'une lecture incertaine ; c'est le bon nombre. Le
recensement a été vérifié sur l'icosaèdre régulier.}. Groupe de symétrie
$\mathbf{Z}/2 \times (\mathbf{Z}/2 \times \mathbf{Z}/2)$, d'ordre $8 = 120/15$,
$15 = 30/2$.

La notation $\Pi_{x,y}$ reviendra pages 18 à 20 : $\{\Pi_{s,u},
\Pi_{s,v}\}$ y sera l'image de l'arête $e$ par la permutation $\rho$.

\pagerange{10}{16}
\section*{II. Géométries d'incidence et ensembles à opérateurs (pages 10 à 16)}

\pagerange{10}{12}
\subsection*{Drapeaux, repères, et le théorème de pleine fidélité}

Soit $\Pi = (S, A, F ; R, R')$ une géométrie d'incidence de dimension $2$,
c'est-à-dire, ici, des données comme à la page 2 satisfaisant au moins les
axiomes 1) et 2)\footnote{La page renvoie à « p.~1 » pour l'ensemble ordonné
$\Phi$ ; c'est vraisemblablement la page 2. Elle cite elle-même « axiomes 1),
2) » pour les degrés qui suivent.}. Un \emph{drapeau} est une partie non vide
totalement ordonnée de $\Phi$, de \emph{longueur} son cardinal moins un ; la
longueur maximale est $2$, et les drapeaux de longueur $2$ sont les
\emph{repères}
\[
E = \operatorname{Rep}(\Pi) = \{(s, a, f) \in S \times A \times F \mid s < a < f\}.
\]
Les drapeaux de longueur $0$ sont les facettes, ceux de longueur $1$ les
éléments de $R$, $R'$, $R''$. Les trois applications d'oubli
$E \to R'$, $E \to R''$, $E \to R$, qui oublient respectivement le sommet,
l'arête, la face, sont de degré $2$ par les axiomes 1) et 2) ; elles définissent
trois involutions sans point fixe $\sigma_0, \sigma_1, \sigma_2$ de $E$, où
$\sigma_i$ change la composante de dimension $i$ du repère et garde les deux
autres. On a
\[
\sigma_0\sigma_2(s, a, f) = \sigma_2\sigma_0(s, a, f) = (s', a, f'),
\]
$s'$ étant l'autre sommet et $f'$ l'autre face de $a$, de sorte que $\sigma_0$
et $\sigma_2$ commutent et que $\sigma_0\sigma_2$ est aussi sans point fixe.
L'ensemble $E$ devient ainsi un ensemble à opérateurs sous le groupe
$\widehat{\Gamma}$ des conventions\footnote{Le diagramme de la page 11 ajoute,
entre crochets, la construction en deux étages $E = R \times_A R'$
au-dessus de $S \leftarrow R \to A \leftarrow R' \to F$, toutes les flèches
de degré $2$.}.

\textbf{Théorème} (page 12). \emph{Sur les géométries d'incidence qui
satisfont 1) à 3'), le foncteur $\Pi \mapsto \operatorname{Rep}(\Pi)$, des
géométries avec leurs isomorphismes vers les $\widehat{\Gamma}$-ensembles où
$\sigma_0, \sigma_1, \sigma_2, \sigma_0\sigma_2$ opèrent sans point fixe avec
leurs isomorphismes, est pleinement fidèle.} On reconstitue en effet $\Pi$
à partir de $E$ :
\[
E/\langle \sigma_1, \sigma_2 \rangle \xrightarrow{\ \sim\ } S, \qquad
E/\langle \sigma_0, \sigma_2 \rangle \xrightarrow{\ \sim\ } A, \qquad
E/\langle \sigma_0, \sigma_1 \rangle \xrightarrow{\ \sim\ } F .
\]
La flèche médiane est bijective parce que $\langle \sigma_0, \sigma_2 \rangle
\simeq \mathbf{Z}/2 \times \mathbf{Z}/2$ opère librement sur la fibre d'une
arête, qui a quatre éléments. La fibre d'un sommet $s$ est l'ensemble des
drapeaux du polygone $\Phi_{>s}$, et ses orbites sous $\langle \sigma_1,
\sigma_2 \rangle$ sont les composantes connexes de ce polygone : la première
flèche est bijective exactement par 3'), et la dernière par 3). Enfin $R$ et
$R'$ sont les images de $E$ dans $S \times A$ et $A \times F$\footnote{La page
dit que 3) et 3') expriment « respectivement que la dernière et la première
flèches sont bijectives » --- dans cet ordre, ce qui est juste --- et conclut
par une phrase d'une lecture douteuse, « on peut donc exprimer $S, A, F$ à
partir de $R$ ». Elle ne démontre pas la pleine fidélité au-delà de ces
reconstructions : un isomorphisme de $\widehat{\Gamma}$-ensembles induit des
bijections sur les trois quotients qui respectent les incidences, et
l'isomorphisme de géométries obtenu induit le premier, car $E \to S \times A
\times F$ est injective.}.

Ce codage d'un polyèdre par l'action de trois involutions sur ses drapeaux est
celui qu'utilisent aujourd'hui la théorie des cartes et celle des polytopes
abstraits\footnote{Le foncteur n'est pas essentiellement surjectif, et la page
ne le prétend pas : un $\widehat{\Gamma}$-ensemble transitif quelconque est
l'ensemble des drapeaux d'une carte, où deux faces peuvent par exemple se
toucher le long de plusieurs arêtes, et non toujours d'une géométrie
d'incidence.}.

\pagerange{13}{14}
\subsection*{Automorphismes, connexité, régularité}

\textbf{Corollaire.} $\operatorname{Aut}(\Pi) =
\operatorname{Aut}_{\widehat{\Gamma}}(\operatorname{Rep}(\Pi))$.

\textbf{Corollaire.} Pour $\Pi$ non vide, conditions équivalentes : a) deux
sommets sont toujours reliés par une chaîne d'arêtes ; a') deux faces sont
toujours reliées par une chaîne de faces adjacentes le long d'arêtes ; b) $\Pi$
n'est pas somme disjointe de deux géométries non vides ; c) $\widehat{\Gamma}$
opère transitivement sur $E$. On dit alors que $\Pi$ est
\emph{connexe}\footnote{La page écrit en c) « espace homogène sous $G_2$ »,
avec un $G$ gras qui diffère du $\mathfrak{S}_2$ de la même ligne ; c'est le
même groupe $\widehat{\Gamma}$. Les équivalences utilisent 3) et 3').}.

\textbf{Corollaire.} Si $\Pi$ est connexe, $\operatorname{Aut}(\Pi)$ opère
\emph{librement} sur $\operatorname{Rep}(\Pi)$ : un automorphisme d'un espace
homogène qui fixe un point est l'identité. Donc, $\Pi$ étant fini,
$\operatorname{Card}\operatorname{Aut}(\Pi)$ divise
$\operatorname{Card}\operatorname{Rep}(\Pi)$. C'est ce que la lecture du
dossier 81 appelle la rigidité d'un objet repéré.

\textbf{Re-corollaire.} Soit $\Pi$ connexe, $r \in E$, $H \subset
\widehat{\Gamma}$ son stabilisateur et $G = \operatorname{Aut}(\Pi)$.
Conditions équivalentes :
\begin{enumerate}
  \item[a)] $G$ est transitif sur $E$ ;
  \item[b)] $H$ est distingué dans $\widehat{\Gamma}$ ;
  \item[c)] ($\Pi$ fini) $\operatorname{Card} G = \operatorname{Card} E$ ;
  \item[d)] pour tout $g \in \widehat{\Gamma}$, il existe $\bar g \in G$ avec
    $g \cdot r = \bar g \cdot r$ ;
  \item[d')] pour $i = 0, 1, 2$, il existe $\bar\sigma_i \in G$ avec
    $\sigma_i \cdot r = \bar\sigma_i \cdot r$.
\end{enumerate}
Pour $E = \widehat{\Gamma}/H$, le groupe des automorphismes est $N(H)/H$, d'où
a) $\Leftrightarrow$ b) ; c) $\Leftrightarrow$ a) parce que l'action est
libre ; d) n'est que a) puisque $E = \widehat{\Gamma} r$ ; et d') entraîne d)
parce que les automorphismes commutent aux $\sigma_i$ : si $h r = \bar h r$,
alors $\sigma_i h r = \bar h \sigma_i r = \bar h \bar\sigma_i r$. On dit alors
que $\Pi$ est \emph{régulière}, ou que c'est un \emph{polyèdre combinatoire
régulier}.

\textbf{Corollaire.} Si $(\Pi_0, r_0)$ est un polyèdre régulier muni d'un
repère et $\Pi$ un polyèdre isomorphe, $\operatorname{Isom}(\Pi_0, \Pi) \to
\operatorname{Rep}(\Pi)$, $\varphi \mapsto \varphi(r_0)$, est bijective : un
polyèdre régulier repéré est défini à isomorphisme \emph{unique} près.

\textbf{Le groupe.} Dans un polyèdre régulier, $\bar g$ est unique, et
$g \mapsto \bar g$ est un \emph{anti}-homomorphisme surjectif
$\widehat{\Gamma} \to G$ de noyau $H$ : $\overline{hg} = \bar g\,\bar h$, car
$\bar g \bar h r = \bar g h r = h \bar g r = h g r$. Composé avec $x \mapsto
x^{-1}$, c'est un homomorphisme ; ainsi $G \simeq \widehat{\Gamma}/H$, et les
$\bar\sigma_i$ engendrent $G$\footnote{La page, après avoir barré un corollaire
$E/\sigma_0 \simeq \operatorname{Drap}_{12}(\Pi)$, etc., écrit en marge
« $g \mapsto \bar g$ est un hom.\ […] $\to G$ », la suite étant perdue au bas
de la feuille. Le renversement de l'ordre est de l'édition ; il disparaît si
l'on fait opérer $\widehat{\Gamma}$ à droite.}.

\pagerange{15}{16}
\subsection*{Les stabilisateurs d'un polyèdre régulier}

Deux blocs barrés précèdent l'énoncé. Le premier définit une
\emph{orientation} d'une géométrie comme une partie $R^{+} \subset E$ telle que
chaque $\sigma_i$ l'échange avec son complémentaire $R^{-}$ --- de façon
équivalente, $R^{+}$ s'envoie bijectivement sur chacun des trois ensembles de
drapeaux de longueur $1$ ---, ses éléments étant les repères directs. Le
second annonce que le stabilisateur d'un sommet est diédral\footnote{La page
écrit en b) $R^{+} \simeq \prod_{j \neq i} \Phi_j$ ; c'est le produit fibré,
l'ensemble des drapeaux de type complémentaire de $i$. Les deux blocs restent
lisibles sous les diagonales qui les barrent.}.

\textbf{Proposition} (pages 15 et 16, marquée « Exer »). \emph{Soit $\Pi$ un
polyèdre combinatoire régulier connexe, $E = \operatorname{Rep}(\Pi)$, $G =
\operatorname{Aut}(\Pi)$.}
\begin{enumerate}
  \item[a)] \emph{Soit $p$ l'ordre de $\sigma_1\sigma_2$ et $q$ celui de
    $\sigma_0\sigma_1$ sur $E$. Par tout sommet passent exactement $p$ arêtes,
    et toute face a exactement $q$ côtés.} En effet, $\langle \sigma_1, \sigma_2
    \rangle$ opère sur les $2p(s)$ repères de sommet $s$ comme le groupe diédral
    du polygone $\Phi_{>s}$, où $\sigma_1\sigma_2$ est une rotation d'ordre
    $p(s)$ ; la régularité rend $p(s)$ constant. Par le corollaire 3 et son
    dual, $3 \leq p, q$, et $p, q \leq \infty$ si l'on admet des polyèdres
    infinis\footnote{La page écrit « NB $3 \leq p, q \leq +\infty$ ». La borne
    $3$ demande l'axiome 4) ; pour une géométrie d'incidence qui ne satisfait que
    1) à 3'), $p$ ou $q$ peut valoir $2$.}.
  \item[b)] \emph{Pour $s \in S$, l'homomorphisme $G_s \to
    \operatorname{Aut}(\Phi_{\geq s})$ est un isomorphisme ; en particulier
    $G_s \simeq D_p$.}
  \item[b')] \emph{Pour $f \in F$, $G_f \to \operatorname{Aut}(\Phi_{\leq f})$
    est un isomorphisme ; $G_f \simeq D_q$.}
  \item[c)] \emph{Pour $a \in A$, soit $\sigma$ l'ensemble de ses deux sommets
    et $\varphi$ celui de ses deux faces. Alors $G_a \to \mathfrak{S}_\sigma
    \times \mathfrak{S}_\varphi \simeq \mathbf{Z}/2 \times \mathbf{Z}/2$ est un
    isomorphisme.}
\end{enumerate}
Dans les trois cas, le stabilisateur opère librement et transitivement sur les
repères qui contiennent la facette, et le groupe d'arrivée aussi : l'injection
est une bijection\footnote{Ces justifications sont de l'édition. Ce sont,
dans le langage du dossier 88, les ordres $p$ et $q$ des rotations $\rho_s =
\sigma_1\sigma_2$ et $\rho_f = \sigma_0\sigma_1$. Deux lignes biffées commencent
au bas de la page 16 une autre proposition.}.

\pagerange{18}{25}
\section*{III. Les trièdres distingués et le groupe alterné (pages 18 à 25)}

\pagerange{18}{19}
\subsection*{La permutation $\rho$ des directions d'arêtes}

Sous la chemise « Cours », $\Pi$ est l'icosaèdre combinatoire, et
$\iota \in G$ son antipodisme --- l'automorphisme qui envoie chaque sommet sur
l'unique sommet à distance $3$ ---, sans point fixe sur $S$, $A$, $F$. Les
quinze paires $\bar a = \{a, \iota a\}$ sont les \emph{directions d'arêtes} :
$\bar A = A/\iota$.

Soit $a = \{s, t\}$ une arête, et $u, v$ les troisièmes sommets des deux faces
qui la bordent. Les cinq voisins de $s$ forment un pentagone
$t, u, \Pi_{s,u}, \Pi_{s,v}, v$, où $\Pi_{s,u}$ est le voisin de $s$ adjacent
à $u$ autre que $t$ (comme page 8) ; le côté de ce pentagone opposé à $t$ est
l'arête $\{\Pi_{s,u}, \Pi_{s,v}\}$. Partant de $t$ au lieu de $s$, on trouve
$\{\Pi_{t,u}, \Pi_{t,v}\}$, qui en est l'antipode. D'où une application
$\rho_0 : A \to \bar A$, compatible avec les automorphismes, en particulier
avec $\iota$, qui opère trivialement sur $\bar A$ ; elle passe donc au quotient
en
\[
\rho : \bar A \longrightarrow \bar A .
\]
On calcule $\rho^2(\bar a) = \overline{\{u, \iota v\}} = \overline{\{v, \iota
u\}}$, puis $\rho^3(\bar a) = \bar a$\footnote{Le calcul de $\rho^2$ et de
$\rho^3$ n'est pas détaillé sur la page ; on l'a vérifié sur l'icosaèdre
régulier. La page écrit $\{u, v'\}$ pour $\{u, \iota v\}$. Les indices des
$\Pi$ de la figure marginale sont serrés et parfois surchargés.}. Donc
$\rho^3 = \mathrm{id}$, $\rho$ est une permutation sans point fixe, toutes ses
orbites ont trois éléments, et il y en a cinq. Une orbite s'appelle un
\emph{trièdre distingué}\footnote{« Trièdre » est la lecture retenue, dans les
deux lots de la transcription, d'un mot écrit sans accent visible, en cinq ou
six lettres, et d'abord lu « tridron » ; elle est sûre aux pages 23, 25 et 26.} ;
deux directions sont \emph{orthogonales} si elles sont distinctes et dans le
même trièdre, c'est-à-dire si $\beta = \rho\alpha$ ou $\alpha = \rho\beta$.
Chaque direction a exactement deux directions orthogonales ; la relation
d'orthogonalité est symétrique, la relation $\beta = \rho\alpha$ ne l'est pas.
On note $T$ l'ensemble des cinq trièdres, $\operatorname{tr} : \bar A \to T$
la projection, et l'on voit $T$ aussi comme une partition de $A$ en cinq parties
de six arêtes, stables par $\iota$.

Dans l'icosaèdre régulier, $\{\Pi_{s,u}, \Pi_{s,v}\}$ est orthogonale à
$\{s, t\}$ au sens euclidien, et un trièdre distingué est un repère orthogonal
de l'espace formé de trois axes d'arêtes ; ce sont les axes des cinq octaèdres
dont les sommets sont les milieux de six arêtes de l'icosaèdre, ou encore les axes de coordonnées des cinq façons
d'écrire ses sommets $(0, \pm 1, \pm \phi)$ et permutations circulaires, $\phi$
étant le nombre d'or\footnote{Cette interprétation métrique est de l'édition ;
la page définit l'orthogonalité combinatoirement et ne dit pas qu'elle
coïncide avec l'orthogonalité euclidienne. Une note marginale, « terminologie :
arêtes orthogonales … », est en partie illisible.}.

\pagerange{20}{21}
\subsection*{Les cinq arêtes d'un sommet}

\textbf{Proposition.} \emph{Deux arêtes orthogonales n'ont pas de sommet
commun.} Les quatre arêtes orthogonales à $\{s, t\}$ sont $\{\Pi_{s,u},
\Pi_{s,v}\}$, $\{\Pi_{t,u}, \Pi_{t,v}\}$, $\{u, \iota v\}$ et $\{v, \iota u\}$ ;
aucune ne rencontre $\{s, t\}$.

Ainsi chaque classe de $T$ est un ensemble de six arêtes deux à deux sans
sommet commun, couvrant les douze sommets : la partition $T$ est une
coloration propre des arêtes en cinq couleurs, chaque sommet voyant les cinq
couleurs\footnote{La page 19 s'arrête sur « Interprétation concrète en termes
de coloriage, » au bas de la feuille ; la page 20 commence par la Proposition.
L'interprétation donnée ici --- aujourd'hui, une $1$-factorisation du graphe
de l'icosaèdre --- est de l'édition.}.

\textbf{Corollaire.} \emph{Pour tout $s \in S$, si $P(s)$ est l'ensemble des
cinq arêtes issues de $s$, $a \mapsto \operatorname{tr}(\bar a)$ est une
bijection $P(s) \xrightarrow{\sim} T$} : deux arêtes issues de $s$ ne sont ni
de même direction ni orthogonales, et les deux ensembles ont cinq éléments.

\textbf{Corollaire.} \emph{Le groupe $G_s^{+} \simeq \mathbf{Z}/5$ des
rotations fixant $s$ --- les automorphismes fixant $s$ qui induisent une
rotation du pentagone $\Phi_{\geq s}$ --- fait de $T$ un torseur ; un générateur
opère sur $T$ par une permutation circulaire, et $G_s^{+}$ opère par
permutations paires}, un cycle de longueur impaire étant pair.

\textbf{Le cas de tout $G_s$} (pages 20 et 21). Plus généralement, les
automorphismes d'un polygone combinatoire à $n \equiv 1 \pmod 4$ côtés
induisent des permutations paires de ses sommets et de ses côtés : pour $n$
impair, une rotation est une puissance d'un $n$-cycle, et une symétrie fixe un
sommet et un côté et échange les autres par $(n-1)/2$ transpositions, nombre
pair exactement quand $n \equiv 1 \pmod 4$. Par conséquent, pour un tel
polygone, toutes les numérotations de ses sommets obtenues en le parcourant à
partir d'un repère définissent \emph{la même} orientation de l'ensemble de ses
sommets : un pentagone oriente canoniquement ses cinq sommets\footnote{La page
21 écrit « type $n \equiv 1$ (4) », « canoniquement » et « via épinglage »,
lectures en partie incertaines ; l'épinglage est le choix d'un repère du
polygone. La justification par le décompte des transpositions est de
l'édition.}. Appliqué au pentagone $\Phi_{\geq s}$ et transporté par
$P(s) \simeq T$, cela munit $T$ d'une structure pentagonale $\pi_s$, invariante
par $G_s \simeq D_5$, et d'une orientation $\omega_s$ ; $G_s$ opère fidèlement
sur $T$, par permutations paires.

\pagerange{21}{22}
\subsection*{L'orientation canonique de $T$}

Fixons un repère $r = (s, a, f)$ ; les automorphismes $\bar\sigma_0,
\bar\sigma_1, \bar\sigma_2$ de la page 14 engendrent $G$. Les deux derniers
fixent $s$, donc opèrent sur $T$ par permutations paires ; $\bar\sigma_0$ fixe
$a$ et $f$ en échangeant les extrémités de $a$, donc fixe le troisième sommet
$w$ de $f$, et opère aussi par une permutation paire\footnote{La figure de la
page 22 marque $\sigma_0$, $\sigma_1$, $\sigma_2$ sur les côtés d'un triangle
fondamental hachuré. La fin de la phrase sur $\sigma_0$ est en partie
illisible ; la page note encore $\sigma_i$ les automorphismes $\bar\sigma_i$.}.

\textbf{Corollaire.} \emph{$G$ opère sur $T$ par permutations paires : on a
un homomorphisme $G \to \operatorname{Alt}_T$.}

\textbf{Corollaire.} \emph{Les orientations $\omega_s$, $s \in S$, sont
toutes égales.} En effet $\omega_{g(s)} = g(\omega_s)$ pour $g \in G$, $G$ est
transitif sur $S$, et une permutation paire conserve chaque orientation. On
obtient une \emph{orientation canonique} $\varpi$ de $T$, conservée par $G$.

\pagerange{23}{25}
\subsection*{Une face, une arête, et l'image de $G$}

\textbf{Proposition} (page 23)\footnote{La page lit la première direction du premier triple
« $\overline{\{u, v\}}$ », avec doute sur les lettres et les indices : c'est
$\overline{\{u, v_1\}}$, car $\{u, v\}$ est un côté de $f$, dont le trièdre
$t_w$ est distinct de $t_\omega$ par b). Le second triple est lu comme on
l'écrit. Les deux triples ont été vérifiés sur l'icosaèdre régulier.}. \emph{Soit $f = \{u, v, w\}$ une face, et
$u_1$, $v_1$, $w_1$ les troisièmes sommets des faces qui bordent
respectivement les côtés $\{v, w\}$, $\{w, u\}$, $\{u, v\}$.}
\begin{enumerate}
  \item[a)] \emph{Les directions $\overline{\{u, v_1\}}$, $\overline{\{v,
    w_1\}}$, $\overline{\{w, u_1\}}$ forment un trièdre distingué
    $t_\omega$, et $\overline{\{u, w_1\}}$, $\overline{\{v, u_1\}}$,
    $\overline{\{w, v_1\}}$ un autre, $t_{\omega'}$ ; la rotation d'ordre $3$
    autour de $f$ les permute circulairement. Ils correspondent aux deux
    ordres circulaires $\omega = (u\,v\,w)$ et $\omega' = (u\,w\,v)$ de $f$ :
    $t_\omega$ est formé des directions $\overline{\{x, \omega(x)_1\}}$.}
  \item[b)] \emph{Les trièdres $t_u = \operatorname{tr}\overline{\{v, w\}}$,
    $t_v = \operatorname{tr}\overline{\{w, u\}}$, $t_w =
    \operatorname{tr}\overline{\{u, v\}}$, $t_\omega$ et $t_{\omega'}$ sont
    distincts.}
\end{enumerate}
Pour b), $t_u, t_v, t_w$ sont distincts par le corollaire de la page 20
appliqué aux sommets de $f$, et $t_u \neq t_\omega = \operatorname{tr}
\overline{\{w, u_1\}}$, $t_u \neq t_{\omega'} = \operatorname{tr}\overline{\{v,
u_1\}}$ par le même corollaire appliqué à $w$ et à $v$ ; de même pour $t_v$,
$t_w$.

\textbf{Corollaire.} \emph{On a une bijection canonique}
\[
T \simeq f \sqcup \operatorname{Or}(f),
\]
\emph{qui envoie un sommet $x$ de $f$ sur le trièdre du côté opposé et un ordre
circulaire $\omega$ sur $t_\omega$.} Elle est compatible à l'action de $G_f
\simeq D_3 = \mathfrak{S}_3$, qui opère sur $f$ tautologiquement et sur
$\operatorname{Or}(f)$ par la signature. Donc $G_f$ opère fidèlement sur $T$ ---
une transposition de $f$ y devient un produit de deux transpositions, ce qui
redonne la parité --- et son image dans $\operatorname{Alt}_T$ est un
sous-groupe diédral d'ordre $6$. L'image de $G$ dans $\operatorname{Alt}_T$ a
donc un ordre multiple de $6$, et, par $G_s$, de $10$ : multiple de $30$.

\textbf{Proposition} (page 24). \emph{Soit $a = \{s, t\}$ une arête et
$\varphi = \{u, v\}$ les troisièmes sommets des faces qui la bordent. Les cinq
trièdres $\operatorname{tr}\overline{\{x, y\}}$ ($x \in a$, $y \in \varphi$)
et $\operatorname{tr}\bar a$ sont distincts}, d'où une bijection canonique
\[
T \simeq \{\operatorname{tr}\bar a\} \sqcup (a \times \varphi).
\]
C'est le corollaire précédent appliqué à la face $\{s, t, u\}$ : ses côtés
donnent $\operatorname{tr}\overline{\{t,u\}}$, $\operatorname{tr}
\overline{\{s,u\}}$, $\operatorname{tr}\bar a$, et ses deux trièdres
$t_\omega$, $t_{\omega'}$ contiennent $\overline{\{s, v\}}$ et $\overline{\{t,
v\}}$, $v$ étant le troisième sommet de la face voisine le long de
$a$\footnote{La page dit « résulte des deux prop.\ précédentes » ; la
justification est de l'édition.}. Le groupe $G_a \simeq \mathfrak{S}_a \times
\mathfrak{S}_\varphi \simeq (\mathbf{Z}/2)^2$ (page 16) opère fidèlement sur
$a \times \varphi$, donc sur $T$, par des doubles transpositions : l'ordre de
l'image de $G$ est aussi multiple de $4$, donc de $60 =
\operatorname{Card}\operatorname{Alt}_T$.

\textbf{Théorème} (page 25). \emph{L'homomorphisme canonique $G =
\operatorname{Aut}(\Pi) \to \mathfrak{S}_T$, de l'ensemble des trièdres
distingués de l'icosaèdre, a pour image $\operatorname{Alt}_T$ et pour noyau
le sous-groupe $\mathfrak{z}$ d'ordre $2$ engendré par l'antipodisme.} Le
noyau a $120/60 = 2$ éléments et contient $\iota$, qui fixe chaque direction.

Soit $G^{+}$ le sous-groupe d'indice $2$ des automorphismes qui conservent
l'orientation de l'icosaèdre, c'est-à-dire chacune des deux classes $R^{+}$,
$R^{-}$ de repères de la page 15. L'antipodisme est central et renverse
l'orientation ; donc $G = \mathfrak{z} \times G^{+}$, et :

\textbf{Corollaire} (page 26). \emph{$G^{+} \to \operatorname{Alt}_T$ est un
isomorphisme : le groupe des rotations de l'icosaèdre est le groupe alterné
de ses cinq trièdres distingués.}\footnote{La page écrit d'abord
$G^{+} \xrightarrow{\sim} \mathfrak{A}_5$, la lettre étant d'une lecture
douteuse, puis le corollaire sous la forme qu'on donne. Que $\iota$ renverse
l'orientation est affirmé, non démontré ; dans l'icosaèdre réel, c'est
$-\mathrm{id}$, de déterminant $-1$. L'isomorphisme du groupe des rotations de
l'icosaèdre avec $A_5$ est classique --- c'est le point de départ des
\emph{Leçons sur l'icosaèdre} de Klein (1884) ---, et on le voit d'ordinaire
par l'action sur les cinq octaèdres ou les cinq cubes inscrits. Les feuillets
ne citent rien ; ce qui leur est propre est la définition purement
combinatoire des cinq objets, par la permutation $\rho$, et la construction de
l'orientation $\varpi$ sans choix.}

\pagerange{26}{30}
\section*{Reconstruction de l'icosaèdre gauche par l'ensemble orienté $T$ (pages 26 à 30)}

Le titre de la page 26 annonce la reconstruction de l'icosaèdre
\emph{orienté} ; ce qui est fait est celle de l'\emph{icosaèdre gauche}
$\bar\Pi = \Pi/\iota$, quotient de l'icosaèdre par l'antipodisme : $6$ sommets,
$15$ arêtes, $10$ faces, $60$ repères\footnote{« Icosaèdre gauche » est son
nom, page 28. C'est ce qu'on appelle aujourd'hui l'\emph{hémi-icosaèdre}, carte
régulière de type $\{3, 5\}$ sur le plan projectif réel, dont le graphe est le
graphe complet à six sommets et le groupe $A_5$. Aucun nom de ce genre n'est
sur la page. Comme $\iota$ déplace chaque sommet à distance $3$, le quotient
satisfait encore les axiomes 1) à 4) de la page 2, et $G/\mathfrak{z} \simeq
\operatorname{Alt}_T$ opère simplement transitivement sur ses $60$ repères :
il est régulier.}. Dans tout ce qui suit, $(T, \varpi)$ est l'ensemble des cinq
trièdres muni de son orientation canonique, et $G$ opère sur tout à travers
$\operatorname{Alt}_T$.

\pagerange{26}{27}
\subsection*{Les sommets, les arêtes, les faces}

\textbf{Sommets.} Il y a $\frac12 \cdot 4! = 12$ structures pentagonales sur
$T$. Chacune oriente canoniquement $T$ (page 21) ; les permutations impaires
échangent les deux paquets, et $\operatorname{Alt}_T$, qui contient le
stabilisateur $D_5$ de chacune, opère transitivement sur chaque paquet de
$6$. Soit $\operatorname{Pent}^{+}(T, \varpi)$ le paquet associé à $\varpi$.
L'application $s \mapsto \pi_s$ est compatible avec $G$ et ne change pas quand
on remplace $s$ par $\iota s$ ; son image est dans $\operatorname{Pent}^{+}(T,
\varpi)$ par définition de $\varpi$. Elle est surjective par transitivité, donc
bijective, les deux ensembles ayant $6$ éléments :
\[
S/\iota \xrightarrow{\ \sim\ } \operatorname{Pent}^{+}(T, \varpi).
\]

\textbf{Arêtes.} Soit $\mathbb{A}(T)$ l'ensemble des triplets $(t, \gamma,
\alpha)$, où $t \in T$, $\gamma$ est une structure de carré sur $T
\smallsetminus \{t\}$, et $\alpha$ l'une des deux paires de côtés opposés de
$\gamma$ ; il a $5 \cdot \frac12 (4-1)! \cdot 2 = 30$ éléments\footnote{La page
note d'abord $\mathbb{A}'(T)$, avec un exposant lu comme un accent qui disparaît
ensuite, et définit $\alpha$ comme un élément de « $\operatorname{Ar}(\gamma)/$
antipod. », les côtés de $\gamma$ modulo l'antipodie du carré ; une note de la
marge, « $=$ Ens des $2$ diagonales du … », est en partie illisible. Dans le
triple, la lettre $\gamma$ surcharge un $\pi$.}. Si $\varphi$ est l'autre
paire de côtés opposés, chaque sommet du carré est sur un côté de $\alpha$ et un
côté de $\varphi$, d'où une bijection $T \smallsetminus \{t\} \simeq \alpha
\times \varphi$, donc
\[
T \simeq \{t\} \sqcup (\alpha \times \varphi),
\]
comme pour une arête page 24. Cette bijection oriente $T$ : on numérote $t$,
puis $\alpha \times \varphi$ dans l'ordre lexicographique, et changer l'ordre
de $\alpha$ ou celui de $\varphi$ revient à une double transposition. Ainsi
$\mathbb{A}(T)$ se partage en deux paquets de $15$, échangés par les
permutations impaires, sur chacun desquels $\operatorname{Alt}_T$ opère
transitivement. L'application qui envoie une arête $a$ de $\Pi$ sur le
triplet $(\operatorname{tr}\bar a, \gamma, \alpha)$ obtenu par la bijection de
la page 24, $\alpha$ étant la paire de côtés de $a \times \varphi$ le long
desquels l'extrémité dans $a$ est fixe, est compatible avec $G$ et invariante
par $\iota$ ; son image est donc un des deux paquets, que l'on note
$\mathbb{A}^{+}(T, \varpi)$\footnote{La page écrit l'application sans dire comment $\alpha$ est choisi ;
le choix qu'on fait est celui que demande la proposition de la page 28. Que le
paquet obtenu soit « celui de $\varpi$ » plutôt que l'autre dépend de cette
convention et de l'ordre des facteurs de $\alpha \times \varphi$ ; la page le
fixe par la notation, sans vérification.}, et, comme pour les sommets,
\[
A/\iota \xrightarrow{\ \sim\ } \mathbb{A}^{+}(T, \varpi).
\]

\textbf{Faces.} La face $f$ s'envoie sur la paire $\Omega_f = \{t_\omega,
t_{\omega'}\}$ de la page 23, et $\iota f$ sur la même. Comme
$\operatorname{Alt}_T$ est transitif sur les dix parties à deux éléments,
\[
F/\iota \xrightarrow{\ \sim\ } \mathfrak{P}_2(T).
\]
Le complémentaire $T \smallsetminus \Omega_f = \{t_u, t_v, t_w\}$ est l'image
des sommets de $f$ : c'est sous cette forme, par des parties à trois éléments,
que la page 31 décrira les faces.

\pagerange{28}{30}
\subsection*{Les incidences et les repères}

On écrit $\bar s \leftrightarrow \pi$, $\bar a \leftrightarrow (t, \gamma,
\alpha)$, $\bar f \leftrightarrow \Omega$.

\textbf{Proposition} (pages 28 et 29).
\begin{enumerate}
  \item[a)] \emph{$\bar s$ et $\bar a$ sont incidents si et seulement si $\gamma$
    se déduit de $\pi$ en $t$} : si $u, v$ sont les voisins de $t$ dans $\pi$,
    $\tilde u \neq t$ le voisin de $u$ et $\tilde v \neq t$ celui de $v$,
    $\gamma$ est le carré de diagonales $\{u, \tilde u\}$, $\{v, \tilde v\}$,
    de côtés $\{u, v\}$, $\{v, \tilde u\}$, $\{\tilde u, \tilde v\}$,
    $\{\tilde v, u\}$, et $\alpha = \bigl\{\{u, v\}, \{\tilde u, \tilde
    v\}\bigr\}$. \emph{Par suite, les couples sommet-arête incidents
    s'identifient aux couples $(\pi, t)$ : $\vec A(\bar\Pi) \simeq S(\bar\Pi)
    \times T$.}
  \item[b)] \emph{$\bar a$ et $\bar f$ sont incidents si et seulement si
    $t \notin \Omega$ et $\Omega$ est un côté de $\gamma$ qui n'est pas dans
    $\alpha$.} Une arête est ainsi sur deux faces, et une face a trois côtés,
    un par élément $\tau \notin \Omega$. \emph{Les couples arête-face
    incidents s'identifient aux couples $(\Omega, \Omega')$ de parties à deux
    éléments disjointes}, $\Omega'$ étant le côté opposé à $\Omega$ dans
    $\gamma$.
  \item[c)] \emph{$\bar s$ et $\bar f$ sont incidents si et seulement si
    $\Omega$ est une diagonale --- une « anti-arête » --- de $\pi$} ; les
    couples sommet-face incidents sont les structures pentagonales munies
    d'une diagonale.
  \item[d)] \emph{Les repères de $\bar\Pi$ correspondent aux triplets
    $(\pi, \tau, \Omega)$ avec $\pi \in \operatorname{Pent}^{+}(T, \varpi)$,
    $\tau \in T$, $\Omega$ une diagonale de $\pi$ ne contenant pas $\tau$ et
    contenant exactement un voisin de $\tau$} : pour $\pi$ et $\tau$ donnés
    ($6 \cdot 5$ choix), il y a deux telles $\Omega$, d'où les $60$ repères.
    L'arête est $(\tau, \gamma, \alpha)$, avec $\gamma$ et $\alpha$ déduits de
    $\pi$ en $\tau$ comme en a).
\end{enumerate}
Pour b), étant donnés $\Omega$ et $\tau \notin \Omega$, il y a exactement deux
structures de carré sur $T \smallsetminus \{\tau\}$ ayant $\Omega$ et son
complémentaire $\Omega'$ pour côtés opposés --- elles correspondent aux deux
bijections de non-adjacence entre $\Omega$ et $\Omega'$ ---, et les deux
orientations qu'elles définissent, avec pour $\alpha$ l'autre paire de côtés,
sont opposées : une seule donne un élément de $\mathbb{A}^{+}(T,
\varpi)$\footnote{La page écrit en b) « $\bar a$ et $\bar f$ incidents ssi
$t \in T \smallsetminus \Omega$ », condition nécessaire mais non suffisante
(elle laisse six faces pour deux), et la suite de b) est un palimpseste de
quatre ou cinq lignes biffées où l'on distingue « $\Omega$ sont des arêtes
consécutives » et « donc $\alpha$ est déterminé par ». La condition complète
qu'on donne est de l'édition, et a été vérifiée sur l'icosaèdre régulier ;
c'est elle que suppose le NB de la page 29, qu'on suit ici, et le d). Les
crochets de la page, $\bar R \simeq S(\bar\Pi) \times T(\bar\Pi)$ et $R' =
A^{\uparrow}(\bar\Pi) \simeq \{(\Omega, \Omega') \mid \Omega \cap \Omega' =
\varnothing\}$, sont ceux qu'on a mis dans l'énoncé.}.

\textbf{Les repères comme numérotations} (page 30). Un repère $(\pi, \tau,
\Omega)$ définit une numérotation $T = \{t_0, t_1, t_2, t_3, t_4\}$ : $t_0 =
\tau$, $t_1$ l'élément de $\Omega$ adjacent à $\tau$, et l'on parcourt $\pi$
dans ce sens, de sorte que $\Omega = \{t_1, t_3\}$. Elle est compatible avec
$\varpi$, puisque $\pi \in \operatorname{Pent}^{+}(T, \varpi)$ et que tout
parcours d'un pentagone définit son orientation canonique. Réciproquement, une
numérotation compatible avec $\varpi$ donne $\pi = (t_0\,t_1\,t_2\,t_3\,t_4)$,
$\tau = t_0$, $\Omega = \{t_1, t_3\}$. \emph{Les $60$ repères de l'icosaèdre
gauche sont ainsi les $60$ numérotations de $T$ compatibles avec
$\varpi$}\footnote{La bijection est affirmée par la page ; on a vérifié sur
l'icosaèdre régulier que les $60$ repères de $\bar\Pi$ donnent $60$
numérotations distinctes, toutes de la même orientation. La phrase de la page
30 est inachevée (« définissant un élément de $\operatorname{Rep}(\bar\Pi)$
… »).}.

\textbf{Ce qui reste.} La page 30 s'achève sur un programme : expliciter, en
termes de $(T, \varpi)$ et pour $\Pi$ orienté, les ensembles
$\tilde S(\bar\Pi) \simeq S(\Pi)$, $\tilde A(\bar\Pi) \simeq A(\Pi)$,
$\tilde F(\bar\Pi) \simeq F(\Pi)$ --- c'est-à-dire remonter de l'icosaèdre
gauche à l'icosaèdre ---, et décrire directement $\bar\Pi$ comme polyèdre
quotient. Le premier point reçoit, pour les ensembles, la réponse du tableau
de la page 31 ; les incidences de $\Pi$ ne sont pas écrites.

\pagerange{31}{32}
\section*{Le tableau des drapeaux de l'icosaèdre gauche (pages 31 et 32)}

Ces deux pages, d'une autre encre et d'une main plus posée, reprennent le
résultat dans un cadre abstrait : $E$ est un ensemble à cinq éléments, $\omega$
une orientation de $E$\footnote{Ce $E$ n'est pas l'ensemble des repères de la
deuxième partie ; c'est le $T$ des pages 18 à 30. On garde la lettre de la
page, qui ne sert ici qu'à ce tableau.}. Un diagramme range les ensembles de
drapeaux par type, $D^{0} = S$, $D^{1} = A$, $D^{2} = F$, $D^{01} = \vec A$,
$D^{12} = A^{\uparrow}$, $D^{02} = A^{\wedge}$, $D^{012} = A^{\uparrow\to}$,
avec au-dessus $\tilde S$, $\tilde A$, $\tilde F$ ; le tableau en donne les
descriptions, dont on vérifie les cardinaux\footnote{Pour $A^{\uparrow}$ la page écrit $\operatorname{Tr}_3(E)$, qui n'a
que dix éléments, avec le commentaire « triangles pointés », lu avec doute ;
c'est ce dernier qui est juste. Pour $F$, $A^{\uparrow}$ et $\tilde F$, le
symbole $\operatorname{Tr}_3$ surcharge un premier symbole, peut-être
$\mathfrak{P}_3$, et des commentaires biffés disent « parties à trois
éléments ». La page ajoute pour $A$ « $\simeq$ structures carrées sur une partie à $4$ él. », par la bijection qui fait des deux paires d'une partition les diagonales du carré, et pour $A^{\wedge}$ « $\simeq \operatorname{Pent}^{s}_\omega(E)$ », une arête d'un pentagone étant opposée à un sommet. Une face correspond à la
partie à trois éléments $T \smallsetminus \Omega$ des pages 27 et 28 ; une
diagonale marquée d'un pentagone (page 29) équivaut à une arête marquée, celle
qui lui est opposée.} :
\[
\begin{array}{lll}
S \simeq \operatorname{Pent}_\omega(E) & \text{pentagones compatibles avec } \omega & 6 \\
A \simeq \operatorname{Car}(E) & \text{partitions de type } (2,2,1) & 15 \\
F \simeq \operatorname{Tr}_3(E) & \text{parties à trois éléments} & 10 \\
\vec A \simeq \operatorname{Pent}^{s}_\omega(E) & \text{pentagones à sommet marqué} & 30 \\
A^{\uparrow} & \text{triangles pointés} & 30 \\
A^{\wedge} \simeq \operatorname{Pent}^{a}_\omega(E) & \text{pentagones à arête marquée} & 30 \\
A^{\uparrow\to} \simeq \operatorname{Rep}_\omega(E) & \text{ordres totaux compatibles avec } \omega & 60 \\
\tilde S & \text{permutations circulaires compatibles avec } \omega & 12 \\
\tilde A & \text{carrés orientés} & 30 \\
\tilde F & \text{triangles orientés} & 20
\end{array}
\]
 Les trois dernières lignes ont les cardinaux de $S(\Pi)$,
$A(\Pi)$, $F(\Pi)$ : elles décrivent l'icosaèdre lui-même. On peut lire
$\tilde S$, $\tilde A$, $\tilde F$ comme les quotients de l'ensemble des
repères de $\bar\Pi$ par les rotations $\sigma_2\sigma_1$, $\sigma_0\sigma_2$,
$\sigma_1\sigma_0$ --- un sommet de $\Pi$ est un sommet de $\bar\Pi$ muni d'un
sens de rotation autour de lui ---, ce qui fait de $\Pi$ le revêtement
d'orientation de $\bar\Pi$\footnote{Cette lecture est de l'édition ; la page
donne les flèches du diagramme et les descriptions, sans commentaire.}.

\textbf{Les involutions.} Un repère étant un ordre total $(a_1, \ldots, a_5)$
compatible avec $\omega$, la page donne
\[
\begin{aligned}
\sigma_0(a_1, a_2, a_3, a_4, a_5) &= (a_1, a_4, a_5, a_2, a_3), \\
\sigma_1(a_1, a_2, a_3, a_4, a_5) &= (a_2, a_1, a_5, a_4, a_3), \\
\sigma_2(a_1, a_2, a_3, a_4, a_5) &= (a_1, a_5, a_4, a_3, a_2),
\end{aligned}
\]
trois doubles transpositions des places, et, page 32, en composant de droite à
gauche,
\[
\begin{aligned}
\rho_f = \sigma_1\sigma_0 &: (a_1, \ldots, a_5) \mapsto (a_4, a_1, a_3, a_2, a_5), \\
\rho_s = \sigma_2\sigma_1 &: (a_1, \ldots, a_5) \mapsto (a_2, a_3, a_4, a_5, a_1), \\
\sigma_0\sigma_2 &: (a_1, \ldots, a_5) \mapsto (a_1, a_3, a_2, a_5, a_4),
\end{aligned}
\]
d'ordres $3$, $5$ et $2$ : ce sont les relations du groupe $\widehat{\Gamma}$
et le type $(p, q) = (5, 3)$. Les trois permutations engendrent une action
simplement transitive sur les $60$ ordres compatibles avec $\omega$, isomorphe
à l'action des $\sigma_i$ sur les repères de l'icosaèdre
gauche\footnote{La page écrit la troisième ligne « $\sigma = \sigma_0\sigma_1$ »
; la valeur donnée est celle de $\sigma_0\sigma_2$, d'ordre $2$, alors que
$\sigma_0\sigma_1$ est d'ordre $3$, inverse de $\rho_f$. Les lettres $\rho$ sont
lues avec doute. La page ne dit pas quel dictionnaire relie les ordres aux
repères. Avec celui de la page 30 ($\tau = a_1$, $\Omega = \{a_2, a_4\}$,
$\pi = (a_1 \ldots a_5)$), $\sigma_0$ et $\sigma_2$ sont bien ceux de la page,
mais $\sigma_1$ est $(a_1, \ldots, a_5) \mapsto (a_5, a_4, a_3, a_2, a_1)$ ;
celui de la page change aussi la face. Les deux triples satisfont les mêmes
relations et définissent des $\widehat{\Gamma}$-ensembles isomorphes, ce qu'on
a vérifié par machine : celui de la page correspond à un autre dictionnaire.}.
Le dossier s'arrête là ; le reste de la page 32 est blanc.

\end{document}
